\documentclass{report}
\usepackage[utf8]{inputenc}
\usepackage[T1]{fontenc}
\usepackage{geometry}
\usepackage{graphicx}
\usepackage{float}
\usepackage{hyperref}
\usepackage[english,indonesian]{babel}

\title{Optimasi Pemrosesan Template LaTeX Menggunakan Pendekatan Berbasis AI}
\author{Academic Support Agent}
\date{\today}

\begin{document}
\maketitle

\begin{abstract}

\end{abstract}

\chapter{PENDAHULUAN}


\chapter{TINJAUAN PUSTAKA}


% Original content (commented out):
% %%%%%%%%%%%%%%%%%%%%%%%%%%%%%%%%%%%%%%%%%%%%%%%%%%%%%%%%%%%%%%%%%%%%%%%%%%%%%
% Template para reporte de investigación bajo el proyecto interdisciplinario
% Diseñado por: Dr. Roberto Contreras-Masse y Roberto Contreras-Moheno
% Abril 2023
% Todos los derechos reservados
% Licencia de uso: MIT, la cual permite otros usar, copiar, modificar, fusionar, publicar, distribuir,
% sublicenciar y/o vender copias del trabajo siempre y cuando se mencione al autor original y 
% se incluya el aviso de copyright. 
%
% Copyright (c) 2023 Dr. Roberto Contreras-Masse, Roberto Contreras-Moheno

% Permission is hereby granted, free of charge, to any person obtaining a copy
% of this software and associated documentation files (the "Software"), to deal
% in the Software without restriction, including without limitation the rights
% to use, copy, modify, merge, publish, distribute, sublicense, and/or sell
% copies of the Software, and to permit persons to whom the Software is
% furnished to do so, subject to the following conditions:
%
% The above copyright notice and this permission notice shall be included in
% all copies or substantial portions of the Software.
%
% THE SOFTWARE IS PROVIDED "AS IS", WITHOUT WARRANTY OF ANY KIND, EXPRESS OR
% IMPLIED, INCLUDING BUT NOT LIMITED TO THE WARRANTIES OF MERCHANTABILITY,
% FITNESS FOR A PARTICULAR PURPOSE AND NONINFRINGEMENT. IN NO EVENT SHALL THE
% AUTHORS OR COPYRIGHT HOLDERS BE LIABLE FOR ANY CLAIM, DAMAGES OR OTHER
% LIABILITY, WHETHER IN AN ACTION OF CONTRACT, TORT OR OTHERWISE, ARISING FROM,
% OUT OF OR IN CONNECTION WITH THE SOFTWARE OR THE USE OR OTHER DEALINGS IN
% THE SOFTWARE.
%%%%%%%%%%%%%%%%%%%%%%%%%%%%%%%%%%%%%%%%%%%%%%%%%%%%%%%%%%%%%%%%%%%%%%%%%%%%%%%%%%%



\documentclass{report}

% Packages
\usepackage[left=1in, top=0.75in, right=0.75in, bottom=0.75in]{geometry}
\usepackage[spanish, mexico]{babel}
\usepackage{graphicx} % Required for inserting images
\usepackage{xcolor}
\usepackage[absolute]{textpos}
\usepackage{tikz}
\usepackage{apacite}
\usepackage{booktabs}
\usepackage{lipsum} 

% color definition for Maristas
\definecolor{blueMarista}{RGB}{15, 58, 93}
\definecolor{yellowMarista}{RGB}{250, 181, 0}

% Define variables for document
\def\mytitle{Titulo de la investigación}
\def\autorA{Alumno 1}
\def\autorB{Alumno 2}
\def\autorC{Alumno 3}
\def\autorD{Alumno 4}
\def\autorE{Alumno 5}
\def\maestro{Maestro}

\begin{document}

\thispagestyle{empty} % Remove page numbering
\begin{center}
      \begin{tikzpicture}[remember picture,overlay]
        \fill [color=yellowMarista] (current page.north west) rectangle ([yshift=-0.08\paperheight]current page.north east);
        \fill [color=blueMarista] (current page.north west) rectangle ([yshift=-0.06\paperheight]current page.north east);
        \node [anchor=south] at (current page.south) {\includegraphics[]{tmp_images/cover.png}};
      \end{tikzpicture}

    \bigskip
    
    {\Large \textbf{INSTITUTO MÉXICO DE CIUDAD JUÁREZ}}\\
    \vspace{0.15in}
    {\large {PREPARATORIA}}\\
    \vspace{0.5in}
    Reporte del Proyecto Interdisciplinario\\
    \vspace{0.5in}
    {\huge \textbf{\mytitle}}\\
    \vspace{0.5in}
    que presentan:\\
    \vspace{0.15in}
    {\large {\autorA}}\\
    {\large {\autorB}}\\
    {\large {\autorC}}\\
    {\large {\autorD}}\\
    {\large {\autorE}}\\
    \vspace{0.35in}
    bajo la dirección del profesor\\
    {\large {\maestro}}\\
    \vspace{0.75in}
    Ciudad Juárez, Chih. a {\today}
    
\end{center}

\newpage % Start the document on a new page

\pagenumbering{roman}

\tableofcontents
\listoffigures
\listoftables



\chapter*{Resumen}
\pagenumbering{arabic}

Las referencias usan APA y se usa el archivo referencias.bib, ejemplo \cite{patron2011modelo}

Ejemplo de Tabla. Use la etiqueta para referenciar tablas, como la Tabla \ref{tab:example}
\begin{table}[htbp]
  \centering
  \caption{Example Table}
  \label{tab:example}
  \begin{tabular}{ccc}
    \toprule
    \textbf{Column 1} & \textbf{Column 2} & \textbf{Column 3} \\
    \midrule
    Row 1, Column 1 & Row 1, Column 2 & Row 1, Column 3 \\
    Row 2, Column 1 & Row 2, Column 2 & Row 2, Column 3 \\
    Row 3, Column 1 & Row 3, Column 2 & Row 3, Column 3 \\
    \bottomrule
  \end{tabular}
\end{table}

Ejemplo de ecuaciones. Use la etiqueta para referencia, como en la ecuación \ref{eq:eq1}
\begin{equation}
    \frac{\partial}{\partial t} \left( \rho + \frac{1}{2} \sum_{j=1}^n \mu_j \| \mathbf{v}_j \|^2 \right) + \nabla \cdot \sum_{j=1}^n \rho \mathbf{v}_j = 0
    \label{eq:eq1}
\end{equation}

Ejemplo de figuras. Use la etiqueta para referencia, como está en la Figura \ref{fig:example}
\begin{figure}[htbp]
  \centering
  \includegraphics[scale=0.25]{tmp_images/logotipos maristas-01.png}
  \caption{Example Figure}
  \label{fig:example}
\end{figure}


\chapter{Introduction}
\section{Antecedentes}
\section{Justificación}
\section{Preguntas de investigación}
\section{Objetivos}
\subsection{Generales}
\subsection{Particulares}
\section{Hipótesis}
\section{Estructura del documento}

\chapter{Marco Teórico}
\section{Mapa Conceptual}
\section{Concepto macro 1}
\subsection{Sub Concepto macro 1}

\chapter{Metodología}
\section{Metodología de la investigación}
\section{Materiales y herramientas}
\section{Desarrollo de la investigación}

\chapter{Resultados}
\section{Resultados de la investigación}
\section{Discusión}

\chapter{Conclusiones}

\appendix
\chapter{Appendix Title}

\bibliographystyle{apacite}
\bibliography{referencias}
\end{document}


\end{document}